\documentclass[12pt]{article}
\usepackage[utf8]{inputenc}
\usepackage[T1]{fontenc}
\usepackage[margin=1in]{geometry}
\usepackage{setspace}
\usepackage{xcolor}
\usepackage[colorlinks=true, urlcolor=blue!60!black, linkcolor=blue!60!black]{hyperref}
\setlength{\parindent}{0em}
\setlength{\parskip}{0.6em}
\renewcommand{\baselinestretch}{1.0}
\begin{document}
\begin{center}
{\LARGE \textit{Roe v. Wade}, 410 U.S. 113 (1973)\par}
\vspace{0.5em}
{\large Case Brief --- Two-Page Edition\par}
\end{center}

\section{Memory Jogger}
Due process privacy encompasses the abortion decision, governed by a trimester framework (modified in \textit{Casey}, overruled in \textit{Dobbs}) (\href{https://supreme.justia.com/cases/federal/us/410/113/}{410 U.S. 113}).

\section{Detailed Case Facts}
Texas made abortion a felony unless performed on medical advice to save the mother's life. Norma McCorvey (``Jane Roe''), unmarried and pregnant in Dallas, sued District Attorney Henry Wade; a physician facing prosecution and a married couple joined. Roe gave birth before the case ended.

\section{Procedural History}
A three-judge district court declared the statutes unconstitutional but denied an injunction; after argument and reargument, the Supreme Court affirmed in part and reversed in part.

\section{Judicial Votes}
7--2. Blackmun wrote for Burger, Douglas, Brennan, Stewart, Marshall, and Powell; Burger, Douglas, and Stewart concurred; White (joined by Rehnquist) and Rehnquist dissented.

\section{Holding}
Roe's claim was ``capable of repetition, yet evading review.'' A fundamental privacy right in Fourteenth Amendment liberty includes a qualified abortion decision, and a fetus is not a constitutional person. The State may not regulate the first-trimester decision, may regulate for maternal health afterward, and may proscribe abortion after viability except to preserve the mother's life or health.

\section{Analysis of Opinions}
Blackmun's history found restrictive bans relatively recent, and modern medicine made early abortion safe. Building on \textit{Griswold} and \textit{Eisenstadt}, he located privacy in due process liberty, declined to decide when life begins, and balanced the State's growing interests by trimester. Stewart admitted this was substantive due process. White called it ``an exercise of raw judicial power''; Rehnquist urged rational-basis review, invoked \textit{Lochner}, and cited the 36 states and territories restricting abortion in 1868.

\section{Examples: Future Applications}
\textbf{Same-Side:} a hospital-only rule for first-trimester abortions is invalid. \textbf{Opposite-Side:} prosecuting a non-physician is valid. \textbf{Fence-Sitter:} a 24-hour waiting period with a state script in the first trimester.

\section{Critique}
Originalists, and Justice Alito in \textit{Dobbs}, say the right lacks textual and historical basis. John Hart Ely called it bad constitutional law, and Ruth Bader Ginsburg faulted its sweep and its reliance on privacy rather than equality.

\section{Key Quotations}
\begin{quote}
``This right of privacy, whether it be founded in the Fourteenth Amendment's concept of personal liberty and restrictions upon state action, as we feel it is, or, as the District Court determined, in the Ninth Amendment's reservation of rights to the people, is broad enough to encompass a woman's decision whether or not to terminate her pregnancy.''
\end{quote}
\noindent --- Justice Blackmun, Opinion of the Court, 410 U.S. at 153.

\end{document}
